\section{Introduction}

Quantum computing has undergone rapid development over recent years: from first conceptualization in the 1980s \cite{Benioff1980, Feynman1982}, and early proof of principles for hardware in the 2000s \cite{Chuang1998, Jones1999, Leung2000, Vandersypen2001, Haffner2005, Negrevergne2006, Plantenberg2007, Hanneke2009, Monz2011, Devitt2013, Devitt2016, Monz2016}, quantum computers can now be built with hundreds of qubits \cite{Jurcevic2021, Pino2021, Ebadi2021}. While the technology remains in its infancy, the fast progress of quantum hardware and the massive financial investments all over the world have led many to assert that so-called Noisy-Intermediate Scale Quantum (NISQ) devices \cite{preskillQuantumComputingNISQ2018, Brooks2019} could outperform conventional computers in the near future \cite{Arute2019, Zhong2020, Wu2021, Madsen2022}. NISQ devices are near-term quantum computers, with a limited number of qubits, and too few physical qubits to implement robust error correction schemes. Existing NISQ computers have already been shown to outperform conventional computers on a limited set of problems designed specifically to fit quantum computers' capabilities \cite{Arute2019, Zhong2020, Wu2021}. Algorithms running on these restricted devices may require only a small number of qubits, show some degree of noise resilience, and are often cast as hybrid algorithms, where some steps are performed on a quantum device and some on a conventional computer. In particular, the number of operations, or quantum gates, must remain moderate, as the longer it takes to implement them, the more errors are introduced into the quantum state, and the more likely it is to decohere. Due to these restrictions, there are severe limits on the scope of algorithms that can be considered. In particular, well-known quantum algorithms such as Shor's algorithm \cite{shorAlgorithmsQuantumComputation, Vandersypen2001, ChaoYang2007, Lanyon2007, Lucero2012, MartinLopez2012, Markov2013, Amico2019} for factoring prime numbers, or Grover's algorithm \cite{groverFastQuantumMechanical1996, Bennett1997, Cerf2000, Ambainis2004, Ambainis2007, Bernstein2010} for unstructured search problems, are not suitable.

The Variational Quantum Eigensolver (VQE) was originally developed by Peruzzo {\it et al.} \cite{Peruzzo2014}, and its theoretical framework was extended and formalized by McClean \textit{et al.} in Ref. \cite{mccleanTheoryVariationalHybrid2015}. The VQE is among the most promising examples of NISQ algorithms. In its most general description, it aims to compute an upper bound for the ground-state energy of a Hamiltonian, which is generally the first step in computing the energetic properties of molecules and materials. The study of electronic structures is a critical application for quantum chemistry (for instance: \cite{Deglmann2014, WilliamsNoonan2017, Heifetz2020}) and condensed matter physics (for instance: \cite{Continentino2021, VanderVen2020}). The scope of VQE is therefore very wide-ranging, being potentially relevant for drug discovery \cite{Cao2018_DD, Blunt2022}, material science \cite{Lordi2021}, chemical engineering \cite{Cao2019_QC}. Conventional computational chemistry, grounded in nearly a century of research, provides efficient methods to approximate such properties, but it becomes intractably expensive for very accurate calculations on increasingly large systems. This poses challenges in the practical application of such methods. One of the main reasons why computational chemistry methods can lack sufficient accuracy in molecular systems is an inadequate treatment of the correlations between constituent electrons. These interactions between electrons formally require computation that scales exponentially in the size of the system studied (i.e. the total time it takes to implement the computation is an exponential function of the system size), rendering exact quantum chemistry methods in general intractable with conventional computing. This limitation is well studied in the literature addressing simulation of quantum computers on conventional computers, Ref.~\cite{Zhou2020} provides an excellent example. 

This bottleneck is the motivation for investigating methods such as the VQE, with the anticipation that these could one day outperform the conventional computing paradigm for these problems \cite{Boixo2018, mccaskeyQuantumChemistryBenchmark2019}. In 1982, Richard Feynman theorized that simulating quantum systems would be most efficiently done by controlling and manipulating a different quantum system \cite{Feynman1982}. An array of qubits obey the laws of quantum mechanics, the same way an electronic wavefunction does. The superposition principle \cite{Silverman2008, Ballentine2008} of quantum mechanics means that it can be exponentially costly to encode the equivalent information on conventional devices, while it only requires a linearly growing number of qubits. In the context of electronic structure theory \cite{Helgaker2000, Kratzer2019, Li2020_EST} this is the appeal of quantum computing: it offers the possibility to model and manipulate quantum wavefunctions exactly, beyond what is possible with conventional computing. While largely dominated by electronic structure research, the VQE and its extensions have also been applied to several other quantum mechanical problems which face similar scaling issues. These notably include nuclear physics \cite{Miceli2019, DiMatteo2021} and nuclear structure problems \cite{Kiss2022, Romero2022}, high-energy physics \cite{Bauls2020, Bass2021, Bauer2022}, vibrational and vibronic spectroscopy \cite{McArdle2019_vibra, Sawaya2019, Ollitrault2020_vibrational, Jahangiri2020,Ltstedt2021, Sawaya2021}, photochemical reaction properties predictions \cite{Mitarai2020, Omiya2022}, periodic systems \cite{Liu2020, Yoshioka2022, Manrique2020}, resolution of non-linear Schr{\"{o}}dinger equations \cite{Lubasch2020}, and computation of quantum states of a Schwarzschild-de Sitter black hole or Kantowki-Sachs cosmology \cite{Joseph2022}.

The VQE starts with an initialized qubit register. A quantum circuit is then applied to this register to model the physics and entanglement of the electronic wavefunction. Quantum circuit refers to a pre-defined series of quantum operations that will be applied to the qubits \cite{nielsenQuantumComputationQuantum2010}. The number of consequential operations in a circuit is referred to as depth. This circuit is defined by two parts: (1) a structure, given by a set of ordered quantum gates, often referred to as an 'ansatz'; and (2) a set of parameters that dictates the behavior of some of these quantum gates. Once the quantum circuit has been applied to the register, the state of the qubits is designed to model a trial wavefunction. The Hamiltonian of the system studied can be measured with respect to this wavefunction to estimate the energy. The VQE then works by variationally optimizing the parameters of the ansatz in order to minimize this trial energy, constrained to always be higher than the exact ground state energy of the Hamiltonian by virtue of the variational principle \cite{Rayleigh1870, Ritz1908, Arfken1985}. For the VQE to be tractable, it is necessary that the number of quantum operations required to model the wavefunction is sufficiently low, imposing a relatively compact ansatz. The VQE admits wavefunction ans\"atze which cannot be efficiently simulated on conventional computers, indicating a possible advantage over conventional approaches if these quantum circuits are sufficiently accurate trial wavefunctions \cite{Peruzzo2014}. A first demonstration of the potential of these quantum ans\"atze was shown in \cite{Peruzzo2014}, where an ansatz with polynomially-scaling depth in the size of the qubit register was constructed using principles grounded in conventional quantum chemistry (namely, the Unitary Coupled Cluster, which is discussed in more detail in Sec. \ref{sec:UCCA}). Since then, many alternatives have been proposed, with ans{\"{a}}tze scaling as low a linearly \cite{Lee2019} in the size of the qubit register. It must be understood however that a shallower ansatz, i.e. with fewer necessary quantum operations, will in general cover an overall smaller span of the space of all possible wavefunctions, and could result in lower accuracy of the ground state energy.

While the design of the ansatz is core to the VQE and to the theoretical underpinnings of its potential advantage over conventional methods, there are many other parts of the algorithm that directly affect the cost and feasibility of the approach. We provide a discussion for each of these parts in our review, and briefly outline them here. To begin with, once one has chosen a quantum system to study, one must decide on how to construct its Hamiltonian. It is a requirement of the VQE that the Hamiltonian can be written as a sum of individual terms, where the number of terms increases with system size relatively slowly. This is a feature of electronic Hamiltonians, where the Coulomb interaction, a sum of two-body terms, scales polynomially in the system size \cite{Szabo1996}. Often there is considerable flexibility and freedom in the specific choice or how a Hamiltonian is represented mathematically. This step is critical, as it impacts all aspects of the VQE: number of qubits, depth of the ansatz, and the total number of measurements required. Once a Hamiltonian is constructed, one must translate it into operators that can be directly measured on a quantum computer (spin or Pauli operators), which can again impact the depth of the ansatz and number of measurements. The next step is to decide on an ansatz, which must be expressive enough to be able to approximately model the ground state wavefunction, but not so much that it results in circuits that become too deep, or parameterizations that are too complex to efficiently train. Once an ansatz is selected, one must choose an appropriate optimizer, which significantly impacts the convergence rate of the VQE optimization and the overall cost of the algorithm. Measurements on quantum devices are effectively stochastic \cite{nielsenQuantumComputationQuantum2010}, meaning that to obtain the expectation value of an operator, one must repeat measurements enough times to achieve a given level of precision. The number of measurements required is not only defined by the required precision, but also by the number of operators in the Hamiltonian. Adopting efficient measurement strategies is critical for controlling the implementation cost of VQE. Finally, to reduce the possible impact of quantum noise on the accuracy of the result, one can introduce error mitigation strategies into the individual measurements, which again need consideration of their associated computational cost. 

This review holds two main purposes: (1) provide an overview for each of the different components of the VQE algorithm outlined above, along with a suggestion for best practices found in the literature. Because of the large variety of quantum systems the VQE can handle, we propose sets of best practices for two different and fairly general types of systems: \textit{ab initio} molecular systems and spin-lattice models, though it is likely that some of the conclusions can be generalized to other Hamiltonians, including impurity models, condensed matter, vibrational spectroscopy or nuclear structure; (2) identify a list of open research questions regarding the future applicability of the VQE, and about individual aspects of the algorithm presented above. We identify four significant (and inter-related) hurdles, which impact the future viability of this method. These are presented here at a high level and discussed throughout the review.

Firstly, several studies have considered the scaling of the number of measurements as the system size increases \cite{Wecker2015, Elfving2020, Kuhn2019, Gonthier2020}. Despite scaling polynomially \cite{mccleanTheoryVariationalHybrid2015}, the number of measurements required can nonetheless rapidly become too large for the method to be viable \cite{Wecker2015, Elfving2020, Gonthier2020}. As a consequence, many studies have called for developments in efficient measurement schemes (for instance, Ref.~\cite{Gonthier2020}), with approaches to find compact representations of the Hamiltonian and judicious use of joint measurement of commuting observables an ongoing source of developments in this endeavor (see Sec.~\ref{sec:DecomposedInteractions}). 

Secondly, an alternative approach to the measurement problem could lie in the enormous potential for parallelization of VQE. This potential was clearly called for in the original VQE paper of Peruzzo {\it et al.} \cite{Peruzzo2014}, but has received relatively little consideration from the community, in particular with respect to efficient strategies for this parallelization and possible communication overheads (further discussion on these points is found in Sec.~\ref{sec:parallelization}). 

Thirdly, variational quantum algorithms suffer from a common pitfall in the challenge of parameter optimization: the so-called barren plateau problem \cite{Wecker2015, McClean2018, Cerezo2021_BP}. The barren plateau is a result of the phenomenon that gradients of the VQE parameters vanish exponentially with the number of qubits used in the model \cite{McClean2018}, in the overall expressibility of the ansatz \cite{Holmes2016}, degree of entanglement of the wavefunction modeled \cite{OrtizMarrero2020, Patti2021}, or non-locality of the cost function \cite{Cerezo2021_BP, Uvarov2020, Sharma2020}. This of course would prevent tractable optimization in larger systems. Many methods have been proposed to address this problem, but it is currently unclear whether these methods allow one to fully contain the barren plateau problem in the context of VQE (further discussions on this topic are found in Sec.~\ref{sec:barren_plateau}). A related point is that the complexity of the optimization landscape for the VQE remains poorly understood. Bittel and Kliesch~\cite{Bittel2021} characterized specific features of this landscape and demonstrated the risk of optimization being undermined by local minima and non-convex features. The underlying question here is the convergence rate of different optimizers, and the overall scaling of the number of iterations required to reach convergence as a function of the system size and entanglement in the wavefunction. 

Finally, the extent to which VQE is resilient to quantum noise and its ability to be mitigated in a tractable manner remains unclear. Variational algorithms have the advantage of exhibiting some ability to learn away certain types of systematic noise \cite{Enrico2021EvaluatingNoiseResilience,Enrico2020NoiseInducedBreakingOfSymmetries}. While it was also shown that the cost of several error mitigation methods may largely outweigh the benefits \cite{RyujiFundamentalLimits2021}, the question of the impact of quantum noise on VQE results and the predictability of these errors should be studied in much further depth than has been done so far. Overall, despite a number of impediments, the VQE has the potential to be among the first useful methods to be implemented on NISQ devices. However, as the hardware continues to progress, so must the theoretical underpinning of this approach, as well as the efficiency and robustness of the software and algorithms, to guarantee that the potential benefits of methods such as VQE can be tapped as early as possible. 

\paragraph{Structure of the review:} Sec. \ref{sec:overview} presents a more formal definition of the VQE, its potential advantages and limitations are discussed, and our assessment of the state-of-the-art VQE is presented. The following sections provide reviews of each of the components of the VQE pipeline (as also summarized in Fig. \ref{fig:VQE_pipeline}). The first step in the VQE process is to choose a representation for the molecular Hamiltonian (Sec. \ref{sec:Hamiltonian_representation}), this is followed by a description of the methods used to map fermionic operators to spin operators, therefore allowing the Hamiltonian operators to be directly measured on a quantum device (Sec. \ref{sec:Encoding}). The next section covers efficiencies that can be used to reduce the number of measurements required to estimate the expectation value of a Hamiltonian (Sec. \ref{sec:Grouping}). We then turn to detailing the various quantum circuit structures representing wave function ans\"atze that have been proposed as the core component of the VQE (Sec. \ref{sec:Ansatz}). This is followed by a presentation of the most important optimization methods and how they relate to the VQE (Sec. \ref{sec:Optimization}). Finally, we list and discuss the main error mitigation techniques that could improve the overall accuracy of the algorithm (Sec. \ref{sec:error-mit}). We close our review with a discussion of relevant extensions of the VQE for quantum chemistry applications (Sec. \ref{sec:Extensions_of_VQE}).

\paragraph{Additional resources and other reviews:} Nielsen and Chuang \cite{nielsenQuantumComputationQuantum2010} provide an introductory text detailing the fundamentals of quantum computing and quantum information which remains a seminal work on the topic. For an introductory overview of quantum computing and a pragmatic discussion of the current state of the technology, we recommend the review by Whitfield {\it et al.} \cite{Whitfield2022}. For a more compact version of introductory topics in quantum computing, we point readers to the review by McArdle {\it et al.} \cite{mcardleQuantumComputationalChemistry2018}. This latter reference also provides a comprehensive review of quantum computing methods for quantum chemistry and material sciences. Two reviews on the same topic are provided by Bauer {\it et al.} \cite{bauerQuantumAlgorithmsQuantum2020} and of Motta and Rice \cite{Motta2021_review}. For a review focused on general variational quantum algorithms we recommend the work of Cerezo {\it et al.} \cite{cerezoVariationalQuantumAlgorithms2020}, and for the wider domain of NISQ algorithms, we recommend the review of Bharti {\it et al.} \cite{Bharti2021_review}. There are also reviews dedicated to narrower topics relevant to the VQE, in particular, we recommend two reviews of ans{\"{a}}tze, which complement our analysis: that by Fedorov {\it et al.} \cite{Fedorov2021}, and the review of unitary coupled cluster based-ans{\"{a}}tze by Anand {\it et al.} \cite{Anand2021_review}. Finally, for an overview of the fundamentals of quantum chemistry, we recommend the book of Szabo and Ostlund \cite{Szabo1996}.

The VQE field is expanding rapidly, the current review only attempts to survey the literature to the end of May 2022.